% Abstract: max. 250 words, no abbreviations, no formulas (IJCRSEE rule).
\noindent\textbf{Abstract:} Resting-state electroencephalography is an inexpensive, portable candidate for
screening mild cognitive impairment and dementia, but published deep learning models are large, mix
different recording states and are rarely tested outside their own dataset. This study aimed to build a
compact, interpretable classifier of normal cognition, mild cognitive impairment and dementia and to evaluate
it rigorously. Using 1,187 clinical recordings from Chung-Ang University Hospital, we first corrected a timing
error in the event annotations, reducing the median misalignment in paused recordings from 33 to 3 seconds
and making the eyes-closed, eyes-open and light-stimulation state of every second known. We then designed a
state-aware, age-normative fusion network with 143,000 parameters that combines state-conditioned signal
segments, 48 clinical biomarkers expressed relative to healthy people of the same age, and age, through
interpretable gates. An ensemble of three networks reached 72.9\% accuracy and an area under the curve of 0.87
on the official test split, above the best published single model (70.2\%) at a fraction of its size, and
recognised 95.7\% of healthy participants. Under repeated five-fold cross-validation the network reached
59.5\%, matching logistic regression on the same biomarkers and exceeding two deep learning baselines.
Community volunteers were younger and recorded with a different stimulation protocol than patients, which
inflates apparent performance, and the model transferred to an independent cohort with an area under the curve
of 0.70. We recommend confound-controlled, cross-validated evaluation for such screening tools.

\vspace{6pt}
\noindent\textbf{Keywords:} electroencephalography, dementia, mild cognitive impairment, deep learning,
normative modelling, human--AI interaction.
